% ======================================================================
% Ontology of Continua -- Core 1.3
% Cross-K module: crossk_global_landscape.tex
% Global cross-level landscape
% ======================================================================

\subsubsection{Глобальный ландшафт Cross-K}
\label{sec:crossk-global-landscape}

Ландшафт Cross-K описывает, как отдельные континуумы $K_d$ встроены в единую
многоуровневую структуру. Вместо обработки каждого уровня как изолированной
теории онтология континуумов упорядочивает их в виде цепочки переходов с
ростом размерности:
\[
R_d : K_d \rightarrow K_{d+1},
\]
где каждый переход $R_d$ является ограниченным переходом, который:
\begin{itemize}
  \item сохраняет основные аксиомы (непустоту $\Omega$, пороги, потоки,
        циклы),
  \item вводит хотя бы одну новую ось
        $A_{\text{new}} \in M_d \setminus A(K_d)$,
  \item проходит размерностный порог $T_d > \Theta_{d,\text{dim}}$,
  \item порождает строго более богатое пространство допустимых состояний
        $\Omega(K_{d+1}) \supset \Omega(K_d)$.
\end{itemize}

Следовательно, глобальный ландшафт — это не плоский список уровней, а
структурированная последовательность фазовых переходов, в которой каждый
$K_d$ возникает из предшественника и ограничивает его последователей.

% ----------------------------------------------------------------------
\subsubsection{1. Задействованные уровни и их роли}

В грубом приближении цепочка делится на четыре блока:
\begin{enumerate}
  \item \textbf{Физическая полоса} ($K_0$--$K_2$): базовая структура, действие и
        перколяция связности.
  \item \textbf{Материально-биологическая полоса} ($K_3$--$K_5$):
        поля, химия, компартменты и протонейронная динамика.
  \item \textbf{Когнитивно-социальная полоса} ($K_6$--$K_8$):
        смысл, социальные системы и цивилизационные инфраструктуры.
  \item \textbf{Теоретическая полоса} ($K_9$--$K_{11}$):
        теории, метатеории и транс-метатеоретические ограничения.
\end{enumerate}

Каждый блок внутренне непротиворечив, но также связан явными межуровневыми
операторами $R_d$ и общими метапространствами $M_n$.

% ----------------------------------------------------------------------
\subsubsection{2. Общие оси и пороги между уровнями}

На цепочке $K_0$--$K_{11}$ повторно появляются несколько структурных
мотивов:
\begin{itemize}
  \item \textbf{Оси $A(K_d)$}: каждый уровень вводит новые оси, но наследует
        преобразованные версии ранних (например, энергии, градиентов,
        информации, ролей, символов).
  \item \textbf{Пороги $\Theta(K_d)$}: пороги существования, устойчивости,
        критический, умирания и размерностный пороги появляются на всех
        уровнях, адаптируясь к локальному физическому/биологическому/
        социальному/теоретическому контексту.
  \item \textbf{Домены состояний $\Omega(K_d)$}: у каждого континуума есть
        непустая область допустимых состояний с границей
        $\partial\Omega(K_d)$, определяемой нарушенными порогами.
  \item \textbf{Циклы $C(K_d)$}: минимальные стабилизирующие циклы существуют
        на каждом уровне (от тривиальных физических циклов до институциональных
        и метатеоретических циклов).
\end{itemize}

Глобальный ландшафт можно понимать как эволюцию этих структурных мотивов при
росте размерности и сложности.

% ----------------------------------------------------------------------
\subsubsection{3. Межуровневые потоки, циклы и напряжения}

Межуровневая динамика сводится к трем семействам объектов:
\begin{description}
  \item[Потоки.] Межуровневые потоки $J_{d,d+1}$ отображают состояния и
        структуры из $K_d$ в $K_{d+1}$ (например, химические градиенты,
        формирующие $K_4$, нейронные паттерны, индуцирующие $K_6$,
        социальные нормы, поддерживающие $K_8$).
  \item[Циклы.] Некоторые циклы явно охватывают несколько уровней, например,
        петли ``цивилизация--теория'' ($K_8$--$K_9$--$K_8$) или
        циклы ``метатеория--практика'' ($K_9$--$K_{10}$--$K_8$).
  \item[Напряжение.] Межуровневое напряжение $T_{d,d+1}$ измеряет
        несоответствие ограничений между $K_d$ и $K_{d+1}$; чрезмерное
        напряжение может предотвратить рождение нового уровня или вызвать
        нисходящий коллапс.
\end{description}

Формально можно записать общий функционал межуровневого напряжения:
\[
T_{d,d+1} = F\bigl(
  A(K_{d+1}) - A(K_d),\,
  P(K_{d+1}) - P(K_d),\,
  J_{d,d+1},\,
  \Theta(K_{d+1}) - \Theta(K_d)
\bigr),
\]
где $T_{d,d+1} > \Theta_{d,\text{dim}}$ сигнализирует о размерностном
переходе.

% ----------------------------------------------------------------------
\subsubsection{4. Условия рождения и исчезновения в глобальной цепочке}

Рождение нового уровня $K_{d+1}$ определяется тремя условиями:
\begin{enumerate}
  \item \textbf{Наличие новой оси:}
        существует $A_{\text{new}} \in M_d$ такая, что
        $A_{\text{new}} \notin A(K_d)$.
  \item \textbf{Размерностное напряжение:}
        структурное напряжение превышает критический порог,
        $T_d > \Theta_{d,\text{dim}}$, что делает старую конфигурацию
        неустойчивой без новой размерности.
  \item \textbf{Допустимое расширение состояний:}
        новая область удовлетворяет условиям
        $\Omega(K_{d+1}) \neq \varnothing$ и
        $\Omega(K_{d+1}) \supset \Omega(K_d)$.
\end{enumerate}

Уровень $K_d$ может исчезать локально (крах конкретного
континуума) или глобально (разрушение всего класса) при:
\[
\Omega(K_d) = \varnothing
\quad\text{or}\quad
k(K_d) \rightarrow 0
\quad\text{or}\quad
\tau_{\text{cycle}}(K_d) \rightarrow \infty.
\]

Глобальный коллапс нескольких соседних уровней можно понимать как каскад
нарушения пороговых условий на соседних уровнях.

% ----------------------------------------------------------------------
\subsubsection{5. Канонические межуровневые цепочки}

Отдельные межуровневые цепочки особенно важны:

\begin{itemize}
  \item \textbf{Физико-химико-биологическая цепочка:}
        $K_1 \rightarrow K_2 \rightarrow K_3 \rightarrow K_4 \rightarrow K_5$
        (от действия и перколяции до мембран и протонейронной динамики).
  \item \textbf{Биологически-когнитивно-социальная цепочка:}
        $K_4 \rightarrow K_5 \rightarrow K_6 \rightarrow K_7$ (от градиентов
        и спайков к когниции и институтам).
  \item \textbf{Социально-цивилизационно-теоретическая цепочка:}
        $K_7 \rightarrow K_8 \rightarrow K_9 \rightarrow K_{10} \rightarrow K_{11}$
        (от норм и ролей к цивилизационным технологиям и
        мета-рамкам).
\end{itemize}

Детализированные модули Cross-K в техническом атласе уточняют эти
канонические цепочки, давая явные описания общих осей, порогов, потоков,
циклов и условий исчезновения для каждого смежного перехода.

В таком смысле глобальный ландшафт Cross-K является \emph{картой} того,
как все отдельные континуумы $K_d$ укладываются в единую согласованную,
многоуровневую структуру.
